% GENERATED FILE. Edit canonical lessons, then run npm run generate:books.
% canonical-source-hash: fnv1a64:cede307b5e801239
% canonical-lessons: RU-C414-perron, RU-C414-platforma, RU-C414-tablo, RU-C414-otpravleniye, RU-C414-pribytiye

\chapter{On the platform}
\label{ch:ru-on-the-platform}
\hlchaptermodality{414}

\begin{chapteropening}
\textbf{By the end of this chapter:} \emph{I can find my train: the platform, the departures board, departure and arrival.}

Complete the last lesson of chapter 414: I can find my train: the platform, the departures board, departure and arrival.
\end{chapteropening}

\section[\texorpdfstring{\ru{перрон} (perrón)}{перрон (perrón)}]{\ru{перрон} (perrón) — a station platform}

\label{lesson:RU-C414-perron}



\begin{quote}
Before the new one: say the Russian for a rate, then the Russian for a concession.
\end{quote}

\subsection*{You'll want to know: \ru{перрон}}
\textbf{\ru{перрон}} (\emph{perrón}) — \textquotedblleft{}a station platform\textquotedblright{}.

\begin{quote}
\ru{Мы} \ru{ждали} \ru{поезд} \ru{на} \ru{перроне}.
\end{quote}

\begin{quote}
\emph{My zhdáli póyezd na perróne.}
\end{quote}

\begin{quote}
We waited for the train on the platform.
\end{quote}

Masculine. \emph{\ru{На} \ru{перроне}} — on the platform.

\begin{cousinweb}
\emph{\ru{перрон}} was borrowed from French \emph{perron} (EtymDB), a stone landing with steps outside a building, built on \emph{pierre} 'stone'.
\end{cousinweb}

\begin{grammarlens}[title={What you've built}]
One more word for telling, explaining and describing at B1.
\end{grammarlens}

\subsection*{Your turn}
\begin{itemize}
\raggedright
  \item \emph{Say it:} \emph{perrón}
  \item \emph{Say it:} \emph{perrón}, once more
  \item \emph{Say it:} a sentence of your own with \emph{perrón}
  \item \emph{Recall:} say the Russian for a rate, then the Russian for a concession, then say \emph{perrón} again
\end{itemize}

\begin{tcolorbox}[breakable,colback=teal!4,colframe=teal!35!black,title={Before you move on}]
What does \ru{перрон} mean? (\textquotedblleft{}A station platform\textquotedblright{}.) Say it once more.
\end{tcolorbox}
\section[\texorpdfstring{\ru{платформа} (platfórma)}{платформа (platfórma)}]{\ru{платформа} (platfórma) — a platform}

\label{lesson:RU-C414-platforma}



\begin{quote}
Before the new one: say the Russian for a concession, then the Russian for a station platform.
\end{quote}

\subsection*{You'll want to know: \ru{платформа}}
\textbf{\ru{платформа}} (\emph{platfórma}) — \textquotedblleft{}a platform\textquotedblright{}.

\begin{quote}
\ru{Поезд} \ru{приходит} \ru{на} \ru{вторую} \ru{платформу}.
\end{quote}

\begin{quote}
\emph{Póyezd prikhódit na vtorúyu platfórmu.}
\end{quote}

\begin{quote}
The train comes in at platform two.
\end{quote}

Feminine. \emph{\ru{Поезд} \ru{стоит} \ru{у} \ru{третьей} \ru{платформы}} — the train is at platform three.

\begin{cousinweb}
\emph{\ru{платформа}} is usually traced to French \emph{plate-forme}, from \emph{plat} 'flat' and \emph{forme} 'a form' (EtymDB lists the French word); the Russian word is not in EtymDB.
\end{cousinweb}

\begin{grammarlens}[title={What you've built}]
One more word for telling, explaining and describing at B1.
\end{grammarlens}

\subsection*{Your turn}
\begin{itemize}
\raggedright
  \item \emph{Say it:} \emph{platfórma}
  \item \emph{Say it:} \emph{platfórma}, once more
  \item \emph{Say it:} a sentence of your own with \emph{platfórma}
  \item \emph{Recall:} say the Russian for a concession, then the Russian for a station platform, then say \emph{platfórma} again
\end{itemize}

\begin{tcolorbox}[breakable,colback=teal!4,colframe=teal!35!black,title={Before you move on}]
What does \ru{платформа} mean? (\textquotedblleft{}A platform\textquotedblright{}.) Say it once more.
\end{tcolorbox}
\section[\texorpdfstring{\ru{табло} (tabló)}{табло (tabló)}]{\ru{табло} (tabló) — a departures board, a display board}

\label{lesson:RU-C414-tablo}



\begin{quote}
Before the new one: say the Russian for a station platform, then the Russian for a platform.
\end{quote}

\subsection*{You'll want to know: \ru{табло}}
\textbf{\ru{табло}} (\emph{tabló}) — \textquotedblleft{}a departures board, a display board\textquotedblright{}.

\begin{quote}
\ru{На} \ru{табло} \ru{уже} \ru{есть} \ru{наш} \ru{поезд}.
\end{quote}

\begin{quote}
\emph{Na tabló uzhé yest nash póyezd.}
\end{quote}

\begin{quote}
Our train is already on the board.
\end{quote}

Neuter and indeclinable: \emph{\ru{на} \ru{табло}}.

\begin{cousinweb}
\emph{\ru{табло}} was borrowed from French \emph{tableau} 'a board, a picture' (EtymDB).
\end{cousinweb}

\begin{grammarlens}[title={What you've built}]
One more word for telling, explaining and describing at B1.
\end{grammarlens}

\subsection*{Your turn}
\begin{itemize}
\raggedright
  \item \emph{Say it:} \emph{tabló}
  \item \emph{Say it:} \emph{tabló}, once more
  \item \emph{Say it:} a sentence of your own with \emph{tabló}
  \item \emph{Recall:} say the Russian for a station platform, then the Russian for a platform, then say \emph{tabló} again
\end{itemize}

\begin{tcolorbox}[breakable,colback=teal!4,colframe=teal!35!black,title={Before you move on}]
What does \ru{табло} mean? (\textquotedblleft{}A departures board, a display board\textquotedblright{}.) Say it once more.
\end{tcolorbox}
\section[\texorpdfstring{\ru{отправление} (otpravléniye)}{отправление (otpravléniye)}]{\ru{отправление} (otpravléniye) — departure}

\label{lesson:RU-C414-otpravleniye}



\begin{quote}
Before the new one: say the Russian for a platform, then the Russian for a departures board.
\end{quote}

\subsection*{You'll want to know: \ru{отправление}}
\textbf{\ru{отправление}} (\emph{otpravléniye}) — \textquotedblleft{}departure\textquotedblright{}.

\begin{quote}
\ru{До} \ru{отправления} \ru{поезда} \ru{осталось} \ru{пять} \ru{минут}.
\end{quote}

\begin{quote}
\emph{Do otpravléniya póyezda ostálos pyat minút.}
\end{quote}

\begin{quote}
There are five minutes left before the train departs.
\end{quote}

Neuter; formal, used on boards and in announcements. \emph{\ru{Почтовое} \ru{отправление}} — an item of mail.

\begin{cousinweb}
\emph{\ru{отправление}} is built on \emph{\ru{отправлять}} 'to send off' with the ending \emph{-\ru{ение}}; EtymDB lists it with the sense 'departure'.
\end{cousinweb}

\begin{grammarlens}[title={What you've built}]
One more word for telling, explaining and describing at B1.
\end{grammarlens}

\subsection*{Your turn}
\begin{itemize}
\raggedright
  \item \emph{Say it:} \emph{otpravléniye}
  \item \emph{Say it:} \emph{otpravléniye}, once more
  \item \emph{Say it:} a sentence of your own with \emph{otpravléniye}
  \item \emph{Recall:} say the Russian for a platform, then the Russian for a departures board, then say \emph{otpravléniye} again
\end{itemize}

\begin{tcolorbox}[breakable,colback=teal!4,colframe=teal!35!black,title={Before you move on}]
What does \ru{отправление} mean? (\textquotedblleft{}Departure\textquotedblright{}.) Say it once more.
\end{tcolorbox}
\section[\texorpdfstring{\ru{прибытие} (pribýtiye)}{прибытие (pribýtiye)}]{\ru{прибытие} (pribýtiye) — arrival}

\label{lesson:RU-C414-pribytiye}



\begin{quote}
Before the new one: say the Russian for a departures board, then the Russian for departure.
\end{quote}

\subsection*{You'll want to know: \ru{прибытие}}
\textbf{\ru{прибытие}} (\emph{pribýtiye}) — \textquotedblleft{}arrival\textquotedblright{}.

\begin{quote}
\ru{Время} \ru{прибытия} — \ru{десять} \ru{часов}.
\end{quote}

\begin{quote}
\emph{Vrémya pribýtiya — désyat chasóv.}
\end{quote}

\begin{quote}
The arrival time is ten o'clock.
\end{quote}

Neuter; formal. In speech people say \emph{\ru{приезд}}.

\begin{cousinweb}
\emph{\ru{прибытие}} is built on \emph{\ru{прибыть}} 'to arrive', from \emph{\ru{при}-} 'to, near' and \emph{\ru{быть}} 'to be'. EtymDB lists \emph{\ru{прибыть}} without a deeper source.
\end{cousinweb}

\begin{grammarlens}[title={What you've built}]
One more word for telling, explaining and describing at B1.
\end{grammarlens}

\subsection*{Your turn}
\begin{itemize}
\raggedright
  \item \emph{Say it:} \emph{pribýtiye}
  \item \emph{Say it:} \emph{pribýtiye}, once more
  \item \emph{Say it:} a sentence of your own with \emph{pribýtiye}
  \item \emph{Recall:} say the Russian for a departures board, then the Russian for departure, then say \emph{pribýtiye} again
\end{itemize}

\begin{tcolorbox}[breakable,colback=teal!4,colframe=teal!35!black,title={Before you move on}]
What does \ru{прибытие} mean? (\textquotedblleft{}Arrival\textquotedblright{}.) Say it once more.
\end{tcolorbox}
